\documentclass[11pt]{article}

\usepackage[margin=1.05in]{geometry}
\usepackage{amsmath}
\usepackage{amssymb}
\usepackage{booktabs}
\usepackage{microtype}
\usepackage[T1]{fontenc}
\usepackage{lmodern}
\usepackage[hidelinks]{hyperref}
\usepackage[round]{natbib}

\title{\bfseries Calibrated Forecasting and Match Leverage in the Serie A\\[2pt]
\large A correction and extension of \emph{Soccer Match Prediction in the Serie A}}
\author{Yannik Pitcan\\
\small \texttt{pitcany@gmail.com}}
\date{August 2026}

\begin{document}
\maketitle

\begin{abstract}
\noindent
A 2016 study of mine, accepted for publication in the \emph{International Journal
of Computer Applications} in 2018, reported 53.0\% three-way accuracy on Serie A
match outcomes and compared that favourably against betting-market accuracy of
55.3\%. This paper corrects that study and extends it from ten seasons to
nineteen (7{,}220 matches, 2007--08 to 2025--26). Three corrections matter. The
benchmark should have been the outcome base rate, 40.6\%, not uniform guessing
at 33\%. Bookmaker odds were used as model \emph{features} and then treated as
an independent standard of comparison, which inverts the logic of the test. And
with 218 scored matches, the reported spread across five classifiers lay inside
two standard errors, so no ranking among them was supported.

Rebuilt on a proper likelihood with tuned time decay and rolling-origin
validation, the model reaches 53.4\% accuracy---statistically indistinguishable
from the 2016 figure---and loses to a margin-free closing price by $0.0067$
Ranked Probability Score ($95\%$ CI $[0.0046, 0.0088]$). In a logarithmic
opinion pool with the market, the optimal weight on the structural model is
$0.000$, on validation and on test alike. Refitting the same machinery to shots
on target, a pre-expected-goals proxy for chance creation, earns weight $0.35$
against the goals model but again $0.000$ against the market: two signals, each
carrying information the other lacks, both already priced.

The constructive consequence is that the deliverable is not a better outcome
forecast but what is built on a calibrated one. We define \emph{match
leverage}---the change in a club's probability of achieving its season objective
between winning and losing a given fixture---and compute it for ACF Fiorentina
from 1 March 2026. The ordering inverts intuition: an away fixture against a
relegation rival carried $2.3\times$ the leverage of hosting the eventual
champions.
\end{abstract}

\section{Introduction}

Forecasting association-football outcomes is a mature problem with an
uncomfortable property: the most accessible benchmark, the betting market, is
also a very strong one. Published work in the applied-computing literature
frequently reports accuracy figures in the low fifties and presents them as
competitive without establishing what the relevant floor and ceiling are. My own
2016 study \citep{pitcan2018} is an instance of this, and the present paper
begins by correcting it.

Section~\ref{sec:errata} documents the specific errors. Sections~\ref{sec:data}
and~\ref{sec:methods} describe the corpus and the rebuilt methodology.
Section~\ref{sec:results} reports the central negative result and the
shots-on-target replication. Section~\ref{sec:leverage} defines match leverage
and applies it. Section~\ref{sec:limits} states what the work cannot do.

Two framing errors drove the original conclusion, and both are worth stating
before the arithmetic. First, the benchmark: the study compared its accuracy
against uniform guessing at $33\%$. Home advantage alone places the majority
class well above that---over this corpus the home side wins $44.1\%$ of
matches---so roughly seven percentage points of the apparent margin were never a
margin at all. Second, and more seriously, bookmaker odds were supplied to the
classifiers as input features, after which the resulting accuracy was compared
against the accuracy of the bookmaker. A model given the market's own forecast
should match it as a floor; scoring below it indicates the pipeline is
discarding information rather than adding any. The original noted the symptom
(``the betting odds did not impact predictions very much'') and attributed it to
the odds being uninformative.

\section{Errata to the 2016 study}
\label{sec:errata}

The following are recomputed directly from the confusion matrices printed in
\citet{pitcan2018}. Both matrices total 218 matches, against a stated test set
of 320; the difference is early-season fixtures dropped for insufficient match
history, a discrepancy the text does not reconcile.

\paragraph{A mislabelled table.} Tables 4.1 and 4.2 of the original are both
captioned ``Linear SVM''. Table 4.1 has diagonal $80+29+2 = 111$, i.e.\
$50.9\%$, matching the reported Linear SVM figure. Table 4.2 has diagonal
$77+34+3 = 114$, i.e.\ $52.3\%$, which is the reported AdaBoost figure. Table
4.2 is AdaBoost.

\paragraph{Binary classifiers below the trivial baseline.} Away wins constitute
$60/218 = 27.5\%$ of the test set, so the constant rule ``never predict an away
win'' scores $72.5\%$. The reported away-win accuracies are $72.9$ (Gaussian
SVM), $69.7$ (Linear SVM), $68.4$ (AdaBoost) and $45.9$ (logistic regression).
None materially exceeds the constant rule; three fall below it. On the home
side, home wins are $92/218$, so ``never predict a home win'' scores
$126/218 = 57.80\%$. The reported Linear SVM home accuracy is $57.8$. To three
significant figures this is the constant rule, which implies the classifier
predicted no home win at all.

\paragraph{Model ranking within noise.} With $n = 218$, the standard error of an
accuracy near one half is $\sqrt{0.25/218} \approx 0.034$. The reported
multi-class accuracies span $47.2\%$ to $53.0\%$---under two standard errors end
to end---and the $n=5$ versus $n=7$ comparison moves the three classifiers in
inconsistent directions. The conclusion that a particular classifier performed
best is not supported.

\paragraph{An error in the form weighting.} The original defines a
time-weighted form feature as $\tau'(\{x_1,\dots,x_n\}) = \frac{1}{n}\sum_i x_i
e^{-i}$. A weighted average must be normalised by $\sum_i w_i$, not by $n$. With
$n=5$, $\sum_{i=1}^{5} e^{-i} = 0.578$, so dividing by $5$ shrinks the weighted
feature to roughly $1/8.7$ of the scale of its unweighted counterpart---material
for the kernel and regularised methods used, which are not scale-invariant.
Separately, the decay is severe: $63.6\%$ of the weight falls on the most recent
match and the Kish effective sample size is $2.1$ matches, so a nominal
five-match window carries about two matches of information. The decay rate was
fixed by assertion rather than fitted.

\section{Data}
\label{sec:data}

The corpus is nineteen complete Serie A seasons, 2007--08 through 2025--26,
comprising $7{,}220$ matches, obtained from \texttt{football-data.co.uk}. Each
record carries date, clubs, full-time score and outcome, shots, shots on target
and corners for both sides, and closing-line prices. Bet365 opening prices are
present for all seasons; Pinnacle opening and closing prices are present from
2012--13, and the closing line is preferred wherever available.

One structural feature of the corpus deserves note, because it motivates a
modelling choice below: promotion and relegation mean only eight clubs appear in
all nineteen seasons, while many appear once or twice. Any rating scheme must
therefore handle clubs with little or no top-flight history.

A second feature is substantive rather than technical. Home advantage in Serie A
has eroded materially over the period: the home win rate falls from $47.5\%$
(2007--13) to $40.2\%$ (2023--26), and the ratio of home to away goals from
$1.38$ to $1.14$. A model fitted on the full corpus without time weighting would
misstate the current strength of home advantage.

\section{Methods}
\label{sec:methods}

\subsection{Market probabilities}

Decimal odds are a reciprocal transform of probability, and the reciprocals of a
quoted book sum to more than one by the bookmaker's margin. Rather than supply
raw odds as features, we convert them to probabilities and remove the margin by
Shin's insider-trading model \citep{shin1993}, which apportions more of the
margin to longshots than proportional normalisation does and is the
best-performing practical method in the comparison of \citet{strumbelj2014}. For
a $k$-outcome book with raw inverse odds $\pi_i$ and book sum $\Pi = \sum_j
\pi_j$, the implied probabilities solve
\begin{equation}
p_i(z) = \frac{\sqrt{z^2 + 4(1-z)\,\pi_i^2/\Pi} - z}{2(1-z)},
\qquad \text{with } z \text{ chosen so } \sum_i p_i(z) = 1 .
\end{equation}
The row sum is strictly decreasing in $z$, so the root is located by bisection.

\subsection{Structural model}

Scorelines are modelled with the bivariate Poisson construction of
\citet{dixoncoles1997}. Each club $j$ carries an attack rating $\alpha_j$ and a
defence rating $\delta_j$; for a fixture between home club $h$ and away club
$a$,
\begin{equation}
\log \lambda = \gamma + \alpha_h - \delta_a, \qquad
\log \mu = \alpha_a - \delta_h ,
\end{equation}
with $\gamma$ the home advantage. The joint mass on low scores is corrected by
the factor $\tau(x,y)$ of \citet{dixoncoles1997}, which adjusts the four cells
$(0,0)$, $(0,1)$, $(1,0)$ and $(1,1)$ through a dependence parameter $\rho$; the
independent-Poisson product demonstrably misfits there. Matches are weighted
exponentially by age,
\begin{equation}
w(t) = \exp\!\left(-\xi \, \Delta t\right),
\end{equation}
with $\Delta t$ in days before the forecast origin. Attack ratings are
constrained to sum to zero, which removes the one-dimensional translation
invariance of the pair $(\alpha, \delta)$ and renders the parameters
identifiable. Clubs absent from the training window---newly promoted
sides---take the exposure-weighted league mean rather than an arbitrary
constant.

The decay rate $\xi$ is selected on the validation period rather than assumed.
The grid search returns a clean interior optimum at $\xi = 0.002$ per day, a
half-life of $347$ days, appreciably slower than the $0.0065$ of the original
Dixon--Coles paper.

\subsection{A process-based variant}

Goals are a sparse realisation of chance creation, and a rating fitted to them
inherits finishing noise. Expected goals is the standard remedy; no freely
licensed source covers Serie A back to 2007. We therefore fit the identical
machinery to \emph{shots on target}, the pre-expected-goals proxy the corpus
already contains, and convert the predicted shot rates to goal rates with a
single league-wide finishing factor estimated on the same window. The factor is
deliberately not club-specific, since club finishing variation is the noise the
substitution is intended to remove.

\subsection{Evaluation}

Because the outcome is ordered---home win, draw, away win---the primary metric
is the Ranked Probability Score \citep{epstein1969}, which is the standard
choice in football forecasting following \citet{constantinou2012}:
\begin{equation}
\mathrm{RPS} = \frac{1}{r-1} \sum_{i=1}^{r-1}
\left( \sum_{j=1}^{i} (p_j - a_j) \right)^{\!2} ,
\end{equation}
with $r=3$. RPS penalises a confident home call more heavily when the away side
wins than when the match is drawn; accuracy and the Brier score cannot make that
distinction. Log loss and Brier are reported alongside, and accuracy is retained
solely for comparability with the 2016 study.

All models are scored walk-forward: fitted on matches strictly preceding an
origin, used to forecast the following fortnight, then refitted. The original
used five-fold cross-validation on time-ordered data, which permits a fold
containing March to inform a prediction about January. Hyperparameters, the
pooling weight and the calibration temperature are fixed on
2013--14 to 2018--19 ($n = 2{,}280$); the test period, 2019--20 to 2025--26
($n = 2{,}660$), is scored once. Every reported mean carries a bootstrap
interval, and model comparisons use a paired bootstrap over matches, exploiting
the fact that competing models score the same fixtures.

\subsection{The incremental-information test}

The question the 2016 study should have posed is not whether the model is
accurate but whether it carries information the market has not already priced. A
logarithmic opinion pool answers this directly. For market forecast $p^{M}$ and
model forecast $p^{S}$,
\begin{equation}
p \;\propto\; \left(p^{M}\right)^{1-w} \left(p^{S}\right)^{w},
\end{equation}
with $w$ fitted out of sample. A weight indistinguishable from zero is evidence
that the model adds nothing; a weight materially above zero is evidence of
genuine incremental information. The logarithmic pool is preferred to a linear
one as it is externally Bayesian and operates naturally on log-odds.

\section{Results}
\label{sec:results}

\begin{table}[t]
\centering
\caption{Test period, 2019--20 to 2025--26 ($n = 2{,}660$). Lower RPS is better.}
\label{tab:main}
\small
\begin{tabular}{lccccc}
\toprule
Model & RPS & 95\% CI & Log loss & Brier & Accuracy \\
\midrule
Market (Shin de-vigged)      & \textbf{0.1905} & [0.1856, 0.1957] & 0.9620 & 0.5717 & 54.8\% \\
Market $+$ model pool        & 0.1905 & [0.1856, 0.1957] & 0.9620 & 0.5717 & 54.8\% \\
Dixon--Coles (calibrated)    & 0.1972 & [0.1921, 0.2024] & 0.9858 & 0.5862 & 53.4\% \\
Shots on target              & 0.2001 & [0.1959, 0.2042] & 0.9959 & --- & 52.9\% \\
Base rate                    & 0.2318 & [0.2292, 0.2344] & 1.0846 & 0.6574 & 40.6\% \\
Uniform (2016 benchmark)     & 0.2340 & [0.2312, 0.2367] & 1.0986 & 0.6667 & 40.6\% \\
\bottomrule
\end{tabular}
\end{table}

Table~\ref{tab:main} reports the central comparison. The market attains RPS
$0.1905$; the structural model $0.1972$. The paired bootstrap places the
difference at $+0.0067$ RPS with $95\%$ interval $[0.0046, 0.0088]$---small in
magnitude but decisively non-zero. The model is worse than the market, not
comparable to it.

Two secondary observations follow. First, the rebuilt model attains $53.4\%$
accuracy against the 2016 study's $53.0\%$. Nine additional seasons, a
principled likelihood, tuned decay and leak-free validation bought $0.4$
accuracy points, well inside noise. This is informative rather than
disappointing: three-way match accuracy appears to be near its ceiling, and
effort spent improving it is spent on a quantity that has stopped moving.
Second, the correct floor is the base rate at $40.6\%$, not uniform guessing.

\paragraph{The pooling weight.} Fitted on validation, the optimal weight on the
structural model in a logarithmic pool with the market is $0.000$. This is not a
small weight but a boundary solution: log loss increases monotonically as any
weight is transferred onto the model, on the validation period and---were one to
cheat and fit it there---on the test period also.

Widening the scan to negative weights, which fall outside the admissible range
for a pool, the validation minimum is interior at $\hat w = -0.225$; carried to
the test period it improves log loss by $0.00062$ ($95\%$ CI
$[-0.00105, 0.00231]$) and RPS by $0.00021$ ($95\%$ CI $[-0.00029, 0.00071]$),
both intervals containing zero. A negative weight divides the market forecast
by the flatter structural one and so sharpens it, which is a temperature
correction on the price rather than a contribution of information.

\begin{table}[t]
\centering
\caption{Pool weights fitted on the validation period.}
\label{tab:pool}
\small
\begin{tabular}{lccc}
\toprule
Pool & Market & Goals model & Shots model \\
\midrule
Market $+$ shots            & 1.00 & ---  & 0.00 \\
Goals $+$ shots             & ---  & 0.65 & \textbf{0.35} \\
Market $+$ goals $+$ shots  & 1.00 & 0.00 & 0.00 \\
\bottomrule
\end{tabular}
\end{table}

\paragraph{Is the target simply too noisy?} The natural objection is that the
model failed because goals are noisy, not because the method is inadequate.
Table~\ref{tab:pool} addresses this. Against the goals model alone, the
shots-on-target variant earns weight $0.35$: it demonstrably carries information
the goals model lacks, so the negative result is not an artefact of target
choice. Against the market, the same signal earns $0.00$, and in the three-way
pool both structural models collapse to zero simultaneously. Two model families
built on different targets, each holding information the other lacks, and the
closing price has already absorbed both.

\paragraph{Calibration and discrimination are distinct.} On the home-win margin
the structural model is \emph{better calibrated} than the market---calibration
slope $0.995$ against the market's $1.103$, with comparable expected calibration
error ($0.021$ for both)---while being clearly less sharp. The fitted
recalibration temperature is $0.99$, i.e.\ essentially no correction is
warranted. The market's advantage is therefore information, not honesty. These
are different deficiencies and they call for different remedies; a study
reporting only accuracy cannot distinguish them.

\section{Match leverage}
\label{sec:leverage}

If a structural model cannot beat a closing price on match outcomes, the product
worth building is not a match-outcome forecast. A club can read the market for
free. What a club cannot read off the market is the quantity that governs its
decisions: how much a particular fixture matters.

Let $\Omega$ denote a season objective---survival, European qualification, the
title---expressed as a band of final league positions. For a remaining fixture
$m$, define
\begin{equation}
L_m \;=\; \Pr\!\left(\Omega \mid \text{win } m\right) \;-\;
          \Pr\!\left(\Omega \mid \text{lose } m\right).
\end{equation}
We estimate $L_m$ by simulating the remainder of the season $20{,}000$ times,
sampling complete scorelines from each fixture's joint distribution so that goal
difference resolves naturally, and conditioning on the simulated result of $m$.
Because fixtures are sampled independently, conditioning after the fact is
equivalent to intervening on $m$ and resampling the remainder; this equivalence
was verified by exhaustive enumeration on a reduced league and by direct
comparison against forced-outcome simulation at $10^{6}$ replications.

\begin{table}[t]
\centering
\caption{ACF Fiorentina, simulated from 1 March 2026. The club stood 16th on 24
points with twelve fixtures remaining. Objective: survival, i.e.\ a final
position of 17th or better.}
\label{tab:leverage}
\small
\begin{tabular}{lccccc}
\toprule
Fixture & Venue & $\Pr(\Omega \mid W)$ & $\Pr(\Omega \mid D)$ & $\Pr(\Omega \mid L)$ & Leverage \\
\midrule
Lecce      & A & 0.987 & 0.959 & 0.914 & \textbf{0.073} \\
Cremonese  & A & 0.986 & 0.957 & 0.922 & 0.063 \\
Verona     & A & 0.982 & 0.952 & 0.927 & 0.055 \\
Genoa      & H & 0.984 & 0.962 & 0.932 & 0.053 \\
Udinese    & A & 0.985 & 0.961 & 0.933 & 0.052 \\
Sassuolo   & H & 0.982 & 0.956 & 0.932 & 0.050 \\
Lazio      & H & 0.988 & 0.960 & 0.942 & 0.046 \\
Parma      & H & 0.981 & 0.955 & 0.936 & 0.045 \\
Atalanta   & H & 0.990 & 0.966 & 0.949 & 0.041 \\
Juventus   & A & 0.988 & 0.971 & 0.951 & 0.037 \\
Roma       & A & 0.988 & 0.968 & 0.952 & 0.036 \\
Inter      & H & 0.988 & 0.971 & 0.955 & \textbf{0.033} \\
\bottomrule
\end{tabular}
\end{table}

Table~\ref{tab:leverage} reports the result. The ordering inverts intuition. The
fixtures that move survival most are those against direct relegation rivals;
those against the leading clubs move it least, because the model already expects
defeat in both branches and has priced it. Losing at home to Inter---the
eventual champions, and by any measure the season's marquee fixture---costs
Fiorentina $3.3$ percentage points of survival probability. Losing away at Lecce
costs $7.3$. The away fixture carries $2.3\times$ the leverage of the home one,
and the gap is approximately eight Monte Carlo standard errors wide, so the
ordering is not simulation noise.

This is a rotation and load-management argument expressed in probability. The
Inter fixture is where a fatigued first-choice player is cheapest to rest; Lecce
away is where he is most expensive. No match-outcome forecast conveys that on
its own.

As a coherence check rather than a claim of predictive success, the simulation
projected $41.4$ expected final points against an actual $42$, and assigned a
survival probability of $96.4\%$; Fiorentina finished 15th. A single realisation
from a distribution is weak evidence, but it establishes that the season
aggregation is not obviously broken, which is the minimum bar before a leverage
ranking is acted upon.

\section{Limitations}
\label{sec:limits}

\textbf{No player-level information.} The model has no access to injuries,
suspensions, rest days, or transfers. This is the largest gap between this work
and what a club holds internally, and the most likely location of a genuine edge
over a market price.

\textbf{No expected goals.} Shots on target weight a tap-in and a speculative
long-range effort equally. Licensed event data would sharpen the process-based
variant, though the results in Table~\ref{tab:pool} suggest it would refine the
conclusion rather than reverse it.

\textbf{The benchmark is a closing price.} It incorporates late team news the
model never observes. Some portion of the $0.0067$ RPS gap is information
asymmetry rather than modelling deficiency.

\textbf{Leverage is single-objective.} It conditions on one position band at a
time. A club trading survival against a cup run requires a utility function, not
a probability difference.

\textbf{Tiebreakers are simplified} to points, then goal difference, then goals
scored; Serie A resolves equal points on head-to-head record first. Exact ties
affecting the reported band occur in under $0.01\%$ of simulations.

\section{Reproducibility}

All results derive from public data and a versioned Python package. The
implementation carries 126 tests. Four areas---the margin-removal solver, the
likelihood and its identifiability, the walk-forward harness, and the leverage
estimator---were additionally subjected to independent adversarial audit, which
confirmed the absence of look-ahead bias by positive control and identified four
defects, all corrected before the results reported here were generated. Two of
the corrected defects were tests that asserted nothing: one passed against a
deliberately leaky harness, the other checked only that probabilities summed to
one. Code and data pipeline: \url{https://github.com/pitcany/seriea-leverage}.

\begin{thebibliography}{9}

\bibitem[Constantinou and Fenton(2012)]{constantinou2012}
Constantinou, A. and Fenton, N. (2012).
Solving the problem of inadequate scoring rules for assessing probabilistic
football forecast models.
\emph{Journal of Quantitative Analysis in Sports}, 8(1).

\bibitem[Dixon and Coles(1997)]{dixoncoles1997}
Dixon, M. and Coles, S. (1997).
Modelling association football scores and inefficiencies in the football betting
market.
\emph{Journal of the Royal Statistical Society: Series C}, 46(2), 265--280.

\bibitem[Epstein(1969)]{epstein1969}
Epstein, E. S. (1969).
A scoring system for probability forecasts of ranked categories.
\emph{Journal of Applied Meteorology}, 8(6), 985--987.

\bibitem[Gneiting and Raftery(2007)]{gneiting2007}
Gneiting, T. and Raftery, A. E. (2007).
Strictly proper scoring rules, prediction, and estimation.
\emph{Journal of the American Statistical Association}, 102(477), 359--378.

\bibitem[Pitcan(2018)]{pitcan2018}
Pitcan, Y. (2018).
Soccer match prediction in the Serie A.
Accepted for publication, \emph{International Journal of Computer Applications}.

\bibitem[Rue and Salvesen(2000)]{rue2000}
Rue, H. and Salvesen, {\O}. (2000).
Prediction and retrospective analysis of soccer matches in a league.
\emph{Journal of the Royal Statistical Society: Series D}, 49(3), 399--418.

\bibitem[Shin(1993)]{shin1993}
Shin, H. S. (1993).
Measuring the incidence of insider trading in a market for state-contingent
claims.
\emph{The Economic Journal}, 103(420), 1141--1153.

\bibitem[{\v S}trumbelj(2014)]{strumbelj2014}
{\v S}trumbelj, E. (2014).
On determining probability forecasts from betting odds.
\emph{International Journal of Forecasting}, 30(4), 934--943.

\end{thebibliography}

\end{document}
